Fix $r\ge 2$.

- Let $K_{r+1}^r$ denote the complete $r$-uniform hypergraph on $r+1$ vertices (it has $\binom{r+1}{r}=r+1$ edges).
- Note that $K_r^r$ is the complete $r$-uniform hypergraph on $r$ vertices, which consists of exactly one $r$-edge.

Let $\mathrm{ex}_r(n;K_{r+1}^r)$ be the maximum number of $r$-edges in an $r$-uniform hypergraph on $n$ vertices containing no copy of $K_{r+1}^r$.

Conjecture (Erd\H{o}s--Sauer, as stated): Every $r$-uniform hypergraph $G$ on $n$ vertices can be written as the union of at most $\mathrm{ex}_r(n;K_{r+1}^r)$ many copies of $K_r^r$ and $K_{r+1}^r$ such that no two copies share a $K_r^r$.

Since $K_r^r$ is a single edge, the condition "no two copies share a $K_r^r$" means simply that the chosen copies are \emph{edge-disjoint}.
So the conjecture is equivalently about partitioning the edge set of $G$ into edge-disjoint blocks, each block either a single edge or the full set of $r+1$ edges of some $(r+1)$-vertex clique $K_{r+1}^r$.
